
\makeatletter
\@ifundefined{remarks}{%
  \theoremstyle{remark}
  \newtheorem{remarks}[theorem]{Remarks}%
}{}
\makeatother

\section{The case of an arbitrary base prescheme}
\SGARefTarget{sga3:xiii:section:3}

Suppose first that the base is a field \(k\), not necessarily
algebraically closed. Since conditions \SGARefLink{sga3:xiii:theorem:2:1}{2.1} \SGARefLink{sga3:xiii:theorem:2:1:item:i-bis}{(i bis)}, \SGARefLink{sga3:xiii:theorem:2:1:item:iv-bis}{(iv bis)}, \SGARefLink{sga3:xiii:theorem:2:1:item:v}{(v)},
and \SGARefLink{sga3:xiii:theorem:2:1:item:vi}{(vi)} are invariant under extension of the base field, passage to
the algebraic closure \(\overline{k}\) of \(k\) shows that they are
equivalent to one another, and equivalent to the assertion that
\(H_{\overline{k}}\) contains a Cartan subgroup of
\(G_{\overline{k}}\). When \SGARefLink{sga3:xiii:theorem:2:1:item:i}{this condition} is satisfied, it is still
true (with the notation of \SGARefLink{sga3:xiii:corollary:2:3}{2.3}) that \(k(X)\) is a finite separable
extension of \(k(G)\), of degree \(d\) independent of every extension
of the base. If \(U\) is the largest open subset of \(G\) such that
\(\psi\) induces a finite étale morphism
\[
  \psi^{-1}(U)\longrightarrow U,
\]
then the formation of \(U\) commutes with extension of the base field.
If \(\psi\) is birational, then \(U\) is also the largest open subset
of \(G\) such that \(\psi\) induces an isomorphism
\[
  \psi^{-1}(U)\xrightarrow{\sim}U;
\]
and if \(g\in U(k)\), there is one and only one subgroup \(H'\) of
\(G\), conjugate to \(H\) over \(\overline{k}\), such that
\(H'\ni g\).

A point \(g\in G(k)\) is called \emph{regular} if it is regular as an
element of
\[
  G(\overline{k})=G_{\overline{k}}(\overline{k}).
\]
More generally, the construction of \SGARefLink{sga3:xiii:corollary:2:7}{2.7} gives an open subset of \(G\),
whose formation commutes with every extension of the base field,
called the \emph{open subset of regular points} of \(G\). It is also
characterized by the following property: for every algebraically
closed extension \(K\) of \(k\) and every \(g\in G(K)\), the point
\(g\) is regular in \(G_K\) if and only if \(g\in U(K)\). If
\(g\in U(k)\), then \(g\) is contained in one and only one Cartan
subgroup of \(G\), as we shall prove below under more general
hypotheses.

Let \(G\) be a smooth, separated group prescheme of finite type, with
connected fibers over the prescheme \(S\), and consider the functor
\SGARefTarget{sga3:xiii:section:3:display:cartan-functor}
\[
  \mathcal C:(\mathrm{Sch})^\circ_{/S}\longrightarrow(\mathrm{Ens})
\]
defined by
\[
  \mathcal C(S')
  =
  \left\{
    \begin{array}{c}
      \text{Cartan subgroups (XII~3.1) of \(G_{S'}\)}
    \end{array}
  \right\}.
\]
Suppose that \SGARefLink{sga3:xiii:section:3:display:cartan-functor}{this functor} is representable by a prescheme smooth over
\(S\). In Exposé XV we give a necessary and sufficient condition for
this to hold, but we already know that the hypothesis is satisfied if
\(G\) is affine over \(S\) with locally constant reductive rank
(\SGARefLink{sga3:xii:3.3}{XII~3.3}), or, more generally, if \(G\) locally admits a maximal torus
for the fpqc topology (\SGARefLink{sga3:xii:7.1}{XII~7.1} a)); this is the case, for example, if
\(S\) is the spectrum of a field.

Let \(X\) be the Cartan subgroup of the \(\mathcal C\)-group
prescheme \(G_{\mathcal C}\), the ``universal Cartan subgroup'' of
\(G\). Thus, as a prescheme over \(S\), \(X\) represents the functor
\[
  X(S')
  =
  \left\{
    \begin{array}{c}
      (C,g)\ \bigm|\ C\text{ is a Cartan subgroup of }G_{S'},\\
      g\text{ is a section of }C\text{ over }S'
    \end{array}
  \right\}.
\]
Consider the canonical projection morphism \((C,g)\mapsto g\),
\[
  \psi:X\longrightarrow G.
\]
We then have \SGARefLink{sga3:xiii:theorem:3:1}{the following theorem}.

\begin{theorem}
\SGARefTarget{sga3:xiii:theorem:3:1}
Under the preceding hypotheses on \(G\), and with the preceding
notation, let \(U\) be the set of \(g\in G\) such that \(g\) is a
regular element of its fiber \(G_s\). Then \(U\) is open, and it is
also the largest open subset \(U\) of \(G\) such that \(\psi\)
induces an isomorphism
\[
  \psi^{-1}(U)\xrightarrow{\sim}U.
\]
\end{theorem}

Let us first prove that \(U\) is open. By the representability of
\(\mathcal C\) as a prescheme smooth over \(S\), and since its
structure morphism is plainly surjective, it follows at once that
\(G\) locally admits a Cartan subgroup for the étale topology and
that the nilpotent rank of the fibers of \(G\) is locally constant.
The same holds for the dimension of the fibers of \(G\). After
localizing on \(S\), we may therefore suppose that these two integers
are constant, say \(\rho\) and \(n\). Consider the Killing polynomial
\[
  P_G(t)=t^n+c_1t^{n-1}+\cdots+c_n\in A[t],
  \qquad A=\Gamma(G,\mathcal O_G).
\]
The restriction of this polynomial to the fibers \(G_s\) of \(G\),
and in particular to the fibers at the maximal points of \(S\), is
divisible by \((t-1)^\rho\). This says that certain linear
combinations of the \(c_i\), with integer coefficients, vanish on the
fibers \(G_s\). When \(S\) is reduced, as we may assume in order to
prove that \(U\) is open, it follows that these combinations
themselves vanish. Hence the Killing polynomial itself is divisible
by \((t-1)^\rho\), say\footnote{The source prints the undefined
exponent \(n-r-1\). It is corrected to \(n-\rho-1\), as required by
the degree \(n-\rho\) of the second factor.}
\[
  P_G(t)
  =
  (t-1)^\rho
  \bigl(t^{n-\rho}+b_1t^{n-\rho-1}+\cdots+b_{n-\rho}\bigr).
\]
Let \(b\) be the sum of the coefficients
\(b_0=1,b_1,\ldots,b_{n-\rho}\) of the second factor. By \SGARefLink{sga3:xiii:corollary:2:7}{2.7}, applied
to the fibers of \(G\), we have
\[
  U=G_b,
\]
which proves that \(U\) is open.

To prove that \(\psi^{-1}(U)\to U\) is an isomorphism, SGA~1,
I~5.7 reduces us to checking this fiber by fiber. We are thereby
reduced to the case of a base field, which may be assumed
algebraically closed. There is then a Cartan subgroup \(C\) of \(G\).
If \(N\) is its normalizer, the conjugacy theorem \SGARefLink{sga3:xii:7.1}{XII~7.1} a) and b)
identifies \(\mathcal C\) with \(G/N\), and the morphism
\(\psi:X\to G\) considered here is precisely the one defined in
\SGARefLink{sga3:xiii:section:2}{Section~2}. The conclusion follows from \SGARefLink{sga3:xiii:theorem:2:6:item:b}{2.6 b)}. The same argument also
shows that \(U\) is the largest open subset of \(G\) for which
\(\psi\) induces an isomorphism \(\psi^{-1}(U)\to U\).

\begin{corollary}
\SGARefTarget{sga3:xiii:corollary:3:2}
Under the hypotheses of \SGARefLink{sga3:xiii:theorem:3:1}{3.1}, let \(g\) be a regular section of \(G\),
that is, suppose that \(g(s)\) is a regular point of \(G_s\) for every
\(s\in S\). Then there is one and only one Cartan subgroup \(C\) of
\(G\) such that \(g\) is a section of \(C\).
\end{corollary}

Indeed, the hypothesis on \(g\) says that \(g\) is a section of
\(U\), while the conclusion says that it admits a unique lift to a
section of \(X\). This is simply another way of saying that
\(\psi^{-1}(U)\to U\) is an isomorphism.

Observe now that the open subset \(\psi^{-1}(U)\) of the Cartan
subgroup \(X\) of \(G_{\mathcal C}\) is precisely the open subset of
\(X\) consisting of the points regular in \(G_{\mathcal C}\), by
which we mean regular in their fibers. We thus obtain a natural
``fibration'' of the dense open subset \(U\) of regular points of
\(G\) over the prescheme \(\mathcal C\); its fibers are dense open
subsets of Cartan subgroups of the fibers of \(G_{\mathcal C}\),
namely their open subsets of points regular in \(G\). For example, we
obtain \SGARefLink{sga3:xiii:corollary:3:3}{the following result}, which will be considerably sharpened in
the next Exposé.

\begin{corollary}
\SGARefTarget{sga3:xiii:corollary:3:3}
Let \(G\) be a smooth connected algebraic group over the field \(k\),
and let \(\mathcal T\) be the scheme of maximal tori of \(G\), which
is isomorphic to the scheme of Cartan subgroups of \(G\). Then the
function field \(k(G)\) of \(G\) is isomorphic to the function field
of a smooth connected affine nilpotent algebraic group \(C\) over the
function field \(k(\mathcal T)\) of \(\mathcal T\), namely
\[
  C=\text{``the generic Cartan subgroup of \(G\)''}.
\]
If \(G\) is affine and has unipotent rank zero, that is, if the Cartan
subgroups of \(G_{\overline k}\)\footnote{The source prints the
undefined base field \(\overline K\); it is corrected to
\(\overline k\).} are tori, then \(k(G)\) is a
unirational extension of \(k(\mathcal T)\).
\end{corollary}

Of course, by the generic Cartan subgroup of \(G\) we mean, by an
abuse of language, the Cartan subgroup of \(G_{k(\mathcal T)}\) that
is the generic fiber of \(X\) over \(\mathcal T\). It remains only to
prove the last assertion of \SGARefLink{sga3:xiii:corollary:3:3}{3.3}, which is contained in \SGARefLink{sga3:xiii:lemma:3:4}{the following
well-known result}, due to Chevalley.

\begin{lemma}
\SGARefTarget{sga3:xiii:lemma:3:4}
Let \(k\) be a field, \(T\) a torus over \(k\), and \(k(T)\) the
field of rational functions on \(T\). Then \(k(T)\) is a
unirational extension of \(k\), that is, it is contained in a purely
transcendental extension of \(k\).
\end{lemma}

Indeed, let \(k'\) be a finite separable extension of \(k\) that
splits \(T\) (\SGARefLink{sga3:x:proposition:1:4}{X~1.4}). Then \(T\otimes_k k'\) is a rational
variety: it contains a dense open subset isomorphic to a dense open
subset of affine \(n\)-space \(\mathbb A^n_{k'}\). Hence
\[
  T'
  =
  \prod_{\operatorname{Spec}(k')/\operatorname{Spec}(k)}
  T_{k'}/\operatorname{Spec}(k')
\]
is a rational variety, since it has a dense open subset isomorphic to
a dense open subset of
\[
  \prod_{\operatorname{Spec}(k')/\operatorname{Spec}(k)}
  \mathbb A^n_{k'},
\]
which is isomorphic to affine \(mn\)-space over \(k\), where
\(m=[k':k]\). Consider the norm homomorphism from \(T'\) to \(T\),
which is defined whenever \(T\) is a commutative group scheme over
\(k\). The composite
\[
  T\longrightarrow T'\longrightarrow T
\]
is the \(m\)-th power map on \(T\), and is therefore dominant.
Consequently \(T'\to T\) is dominant, proving that \(T\) is
unirational.

Return to the hypotheses of \SGARefLink{sga3:xiii:theorem:3:1}{3.1}, but suppose in addition that \(G\)
locally admits a maximal torus for the fpqc topology (\SGARefLink{sga3:xii:7.1}{XII~7.1}). Let
\(Y\) be the maximal torus of the Cartan subgroup \(X\) of
\(G_{\mathcal C}\). The morphism \(\psi:X\to G\) then induces a
morphism \(Y\to G\) whose set-theoretic image consists of the
semisimple elements of the fibers of \(G\) (XII~8). Finally, \SGARefLink{sga3:xiii:theorem:3:1}{3.1}
shows that the restriction of \(\psi\) to the open subset
\(Y^{\mathrm{reg}}\) of regular points of \(Y\) induces a closed
immersion
\[
  Y^{\mathrm{reg}}\longrightarrow U=G^{\mathrm{reg}}.
\]
Making explicit the functorial meaning of
\(Z=Y^{\mathrm{reg}}\), we obtain \SGARefLink{sga3:xiii:corollary:3:5}{the following}.

\begin{corollary}
\SGARefTarget{sga3:xiii:corollary:3:5}
Let \(G\) be a smooth, separated \(S\)-group prescheme of finite
type, with connected fibers over the prescheme \(S\), and suppose
that \(G\) locally admits a maximal torus for the fpqc topology. Let
\[
  Z:(\mathrm{Sch})^\circ_{/S}\longrightarrow(\mathrm{Ens})
\]
be the functor defined by
\[
  Z(S')
  =
  \left\{
    \begin{array}{c}
      \text{regular sections of \(G_{S'}\) over \(S'\) that are}\\
      \text{contained in a maximal torus of \(G_{S'}\)}
    \end{array}
  \right\}.
\]
Then \(Z\) is representable by a closed subprescheme, smooth over
\(S\), of the open subset \(U=G^{\mathrm{reg}}\) of \(G\) introduced
in \SGARefLink{sga3:xiii:theorem:3:1}{3.1}.
\end{corollary}

To conclude, let us record \SGARefLink{sga3:xiii:corollary:3:6}{the following result}, which sharpens the
density theorem \SGARefLink{sga3:xiii:theorem:2:1}{2.1} \SGARefLink{sga3:xiii:theorem:2:1:item:i}{(i)} \(\Rightarrow\) \SGARefLink{sga3:xiii:theorem:2:1:item:vi}{(vi)}.

\begin{corollary}
\SGARefTarget{sga3:xiii:corollary:3:6}
Under the hypotheses of \SGARefLink{sga3:xiii:corollary:3:5}{3.5}, let \(C\) be a Cartan subgroup of \(G\),
and consider the morphism
\[
  \varphi:Z\times C\longrightarrow G
\]
defined by
\[
  \varphi(g,h)=\operatorname{ad}(g)h=ghg^{-1}.
\]
Then \(\varphi\) is dominant.
\end{corollary}

It plainly suffices to prove the assertion fiber by fiber, which
reduces us to the case where \(S\) is the spectrum of an
algebraically closed field \(k\). Let \(T\) be the maximal torus of
\(C\), let \(t_0\in T(k)\) be regular in \(G\), and let
\(c_0\in C(k)\) be regular in \(G\). Consider
\(\varphi^{-1}(\varphi(t_0,c_0))\). Its \(k\)-rational points are the
pairs \((t,c)\), with \(t\in Z(k)\) and \(c\in C(k)\), such that
\[
  \operatorname{ad}(t)c=\operatorname{ad}(t_0)c_0,
  \qquad\text{that is,}\qquad
  c=\operatorname{ad}(t^{-1}t_0)c_0.
\]
They are therefore in bijection with the \(t\in Z(k)\) such that
\(\operatorname{ad}(t^{-1}t_0)c_0\in C\), or equivalently, since
\(c_0\) is regular, such that \(t^{-1}t_0\in N\), where \(N\) is the
normalizer of \(C\); in other words, \(t\in N\). Restricting to those
\(t\in Z(k)\) that belong to \(C(k)\) gives an open and closed part
of this fiber. We have therefore found a connected component of
\(\varphi^{-1}(\varphi(t_0,c_0))\) isomorphic to \(T\). The preceding
set-theoretic argument indeed gives an isomorphism of schemes: one
simply replaces \(k\)-valued points by points with values in an
arbitrary \(k\)-prescheme.

It follows that the generic fiber of \(\varphi\) has dimension at
most \(\dim T\), and therefore
\[
  \dim\operatorname{Im}(\varphi)
  \geq
  \dim(Z\times C)-\dim T
  =
  \dim Z+\dim C-\dim T.
\]
But
\[
  \dim Z
  =
  \dim Y
  =
  \dim\mathcal C+\dim T
  =
  \dim G-\dim C+\dim T.
\]
Thus \(\dim\operatorname{Im}(\varphi)\geq\dim G\), and \(\varphi\)
is dominant.
\par\hfill\textsc{q.e.d.}

\begin{remarks}
\SGARefTarget{sga3:xiii:remark:3:7}
The argument shows in addition that the connected component at
\((t_0,c_0)\) of the fiber
\(\varphi^{-1}(\varphi(t_0,c_0))\) is isomorphic to \(T\). In
particular, it is smooth over \(k\) and has the same dimension as the
generic fiber. It follows that \(\varphi\) is in fact smooth at
\((t_0,c_0)\), something that one should also be able to verify by
calculating the tangent map. Consequently, under the hypotheses of
\SGARefLink{sga3:xiii:corollary:3:6}{3.6}, the induced morphism
\[
  Z\times_S C^{\mathrm{reg}}\longrightarrow G^{\mathrm{reg}},
  \qquad C^{\mathrm{reg}}=C\cap G^{\mathrm{reg}},
\]
is smooth. Likewise, the analogous morphism
\[
  Z\times T^{\mathrm{reg}}\longrightarrow Z,
\]
where \(T\) is a maximal torus of \(G\), is smooth. More generally,
for every smooth connected normal algebraic subgroup \(H\) of \(C\)
containing an element \(c_0\in G(k)\) that is regular in \(G\), the
image of \(Z\times H\to G\) is dense in the image of
\(G\times H\to G\).
\end{remarks}

\section{Lie algebras over a field: rank, regular elements, Cartan subalgebras}
\SGARefTarget{sga3:xiii:section:4}

In the remainder of \SGARefLink{sga3:xiii:expose}{this Exposé} we take up the theory developed by
Chevalley in his book \emph{Théorie des Groupes de Lie III}
(Act.\ Sc.\ Ind.\ 1226, Paris 1955); the methods of schemes allow us
to dispense with the hypothesis of characteristic zero. We begin by
recalling in \SGARefLink{sga3:xiii:section:4}{this section} several well-known notions and results.

Let \(\mathfrak g\) be a Lie algebra over a ring \(k\). For every
\(a\in\mathfrak g\), denote by \(\operatorname{ad}(a)\) the
endomorphism
\[
  \operatorname{ad}(a)\cdot x=[a,x]
\]
of \(\mathfrak g\); it is a derivation of the Lie algebra
\(\mathfrak g\). For every derivation \(D\) of \(\mathfrak g\), the
nil-space of \(D\), that is, the union of the kernels of the iterates
of \(D\), is a Lie subalgebra of \(\mathfrak g\), as follows from the
Leibniz formula
\[
  D^m([x,y])
  =
  \sum_{0\leq p\leq m}
  \binom{m}{p}[D^p x,D^{m-p}y].
\]
We set
\[
  \operatorname{Nil}(a,\mathfrak g)
  =
  \text{nil-space of }\operatorname{ad}(a)
  =
  \bigcup_{m\geq0}\operatorname{Ker}\operatorname{ad}(a)^m.
\]
When no confusion can arise, we write simply
\(\operatorname{Nil}(a)\), and call it the \emph{nil-space of \(a\)
in \(\mathfrak g\)}.

\begin{proposition}
\SGARefTarget{sga3:xiii:proposition:4:1}
For every \(a\in\mathfrak g\), its nil-space
\(\operatorname{Nil}(a,\mathfrak g)\) is a Lie subalgebra of
\(\mathfrak g\), equal to its own normalizer.
\end{proposition}

It remains only to prove that it is its own normalizer. In other
words, every element of
\(\mathfrak g/\operatorname{Nil}(a,\mathfrak g)\) annihilated by the
adjoint representation of \(\operatorname{Nil}(a,\mathfrak g)\) on
this quotient must be zero. This is immediate, since every element
of the quotient annihilated by \(\operatorname{ad}(a)\)\footnote{The
source prints \(\operatorname{Ad}(a)\); because
\(a\in\mathfrak g\), the endomorphism here is
\(\operatorname{ad}(a)\).} is zero.

For the remainder of \SGARefLink{sga3:xiii:section:4}{this section}, suppose that \(k\) is a field and
that \(\mathfrak g\) is finite-dimensional over \(k\). We denote by
\(W(\mathfrak g)\) the scheme over \(k\) defined by \(\mathfrak g\);
its points in a \(k\)-algebra \(A\) are the elements of
\(\mathfrak g\otimes_k A\). If \(a\in\mathfrak g\), the
characteristic polynomial of \(\operatorname{ad}(a)\) is also called
the \emph{characteristic polynomial}, or the \emph{Killing
polynomial}, of \(a\) in \(\mathfrak g\):
\[
  P_{\mathfrak g}(a,t)
  =
  t^n+c_1(a)t^{n-1}+\cdots+c_n(a),
\]
where \(n=\operatorname{rank}_k\mathfrak g\) and \(c_i(a)\in k\).
Taking the same polynomial for
\(a\in\mathfrak g\otimes_k A\), where \(A\) is any \(k\)-algebra, we
see that the \(c_i(a)\) arise from uniquely determined sections
\(c_i\) of the structure sheaf of \(W(\mathfrak g)\), that is, from
elements of the symmetric algebra
\[
  A=\operatorname{Sym}_k(\mathfrak g^\vee),
\]
where \(\mathfrak g^\vee\) is the dual of the \(k\)-module
\(\mathfrak g\). When \(k\) is infinite, the \(c_i\) are determined
by the corresponding polynomial functions \(\mathfrak g\to k\), but
this is no longer true when \(k\) is finite.

Let \(r\) be the largest integer such that the Killing polynomial
\[
  P_{\mathfrak g}(t)
  =
  t^n+c_1t^{n-1}+\cdots+c_n
  \in A[t]
\]
is divisible by \(t^r\); thus
\[
  P_{\mathfrak g}(t)
  =
  t^n+c_1t^{n-1}+\cdots+c_{n-r}t^r,
  \qquad c_{n-r}\neq0.
\]
The integer \(r\) is called the \emph{nilpotent rank} of the Lie
algebra \(\mathfrak g\). It is invariant under extension of the base
field.

\begin{proposition}
\SGARefTarget{sga3:xiii:proposition:4:2}
Let \(r\) be the nilpotent rank of \(\mathfrak g\), and let
\(a\in\mathfrak g\). Then
\[
  \operatorname{rank}_k\operatorname{Nil}(a,\mathfrak g)\geq r,
\]
with equality if and only if
\[
  c_{n-r}(a)\neq0.
\]
In this case \(\operatorname{Nil}(a,\mathfrak g)\) is a nilpotent
Lie algebra. We shall see the converse in \SGARefLink{sga3:xiii:corollary:5:7:part:b}{5.7 b)} when
\(\mathfrak g\) is the Lie algebra of an algebraic group \(G\)
smooth over \(k\).
\end{proposition}

The first assertion is immediate, since by definition
\[
  \operatorname{rank}_k\operatorname{Nil}(a,\mathfrak g)
  =
  \text{the multiplicity of the zero root in }P_{\mathfrak g}(a,t).
\]
Suppose now that \(c_{n-r}(a)\neq0\). We prove that
\(\operatorname{Nil}(a)\) is nilpotent, which is equivalent to saying
that \(\operatorname{ad}_{\operatorname{Nil}(a)}(x)\) is a nilpotent
endomorphism for every \(x\in\operatorname{Nil}(a)\). We may suppose
that \(k\) is algebraically closed. Since
\[
  \operatorname{ad}(a)_{\mathfrak g/\operatorname{Nil}(a)}
\]
is injective, there is a nonempty open subset
\(U\subset W(\operatorname{Nil}(a))\) such that, for every
\(x\in U(k)\), the endomorphism
\(\operatorname{ad}(x)_{\mathfrak g/\operatorname{Nil}(a)}\) is
injective and hence
\(\operatorname{Nil}(x)\subset\operatorname{Nil}(a)\). We may further
suppose that \(U\) is contained in the open subset on which
\(c_{n-r}\) does not vanish; that open subset is nonempty because
\(c_{n-r}(a)\neq0\). Then
\(\operatorname{Nil}(x)\) and \(\operatorname{Nil}(a)\) have the same
dimension, so
\[
  \operatorname{Nil}(x)=\operatorname{Nil}(a).
\]
It follows that
\(\operatorname{ad}(x)_{\operatorname{Nil}(a)}\) is nilpotent for
every \(x\in U(k)\). By the principle of extension of algebraic
identities, this remains true for every
\(x\in\operatorname{Nil}(a)\), and hence \(\operatorname{Nil}(a)\) is
nilpotent.

An element \(a\in\mathfrak g\) is called \emph{regular} if
\(c_{n-r}(a)\neq0\), that is, if
\[
  \operatorname{rank}_k\operatorname{Nil}(a,\mathfrak g)=r.
\]
When \(k\) is infinite, this also says that
\(\operatorname{rank}_k\operatorname{Nil}(a,\mathfrak g)\) is as
small as possible as \(a\) varies in \(\mathfrak g\). In every case,
the notion of a regular element of \(\mathfrak g\) is invariant under
extension of the base field. The set of regular points of
\(W(\mathfrak g)\), meaning those that arise from regular points with
values in a suitable extension of \(k\), is open, since it is
\[
  W(\mathfrak g)_{c_{n-r}},
\]
the set of points at which \(c_{n-r}\) is invertible.

\begin{corollary}
\SGARefTarget{sga3:xiii:corollary:4:3}
Let \(a\) be a regular element of \(\mathfrak g\), and let
\(\mathfrak h\) be a Lie subalgebra of \(\mathfrak g\) containing
\(a\). Then \(\mathfrak h\) is nilpotent if and only if
\[
  \mathfrak h\subset\operatorname{Nil}(a,\mathfrak g).
\]
In particular, \(\operatorname{Nil}(a,\mathfrak g)\) is a maximal
nilpotent subalgebra of \(\mathfrak g\).
\end{corollary}

Since \(\operatorname{Nil}(a)\) is nilpotent, the inclusion
\(\mathfrak h\subset\operatorname{Nil}(a)\) implies that
\(\mathfrak h\) is nilpotent. Conversely, if \(\mathfrak h\) is
nilpotent, it is contained in the nil-space of its element \(a\);
that is, \(\mathfrak h\subset\operatorname{Nil}(a)\).

\begin{proposition}
\SGARefTarget{sga3:xiii:proposition:4:4}
Suppose that \(k\) is infinite. Let \(\mathfrak d\) be a Lie
subalgebra of \(\mathfrak g\). Consider the following conditions:
\begin{enumerate}
\item[(i)]\SGARefTarget{sga3:xiii:proposition:4:4:item:i}
\(\mathfrak d\) is maximal nilpotent and contains a regular element
of \(\mathfrak g\).
\item[(i bis)]\SGARefTarget{sga3:xiii:proposition:4:4:item:i-bis}
\(\mathfrak d=\operatorname{Nil}(a,\mathfrak g)\) for some regular
element \(a\) of \(\mathfrak g\).
\item[(ii)]\SGARefTarget{sga3:xiii:proposition:4:4:item:ii}
\(\mathfrak d\) is nilpotent and
\(\mathfrak d=\operatorname{Nil}(a,\mathfrak g)\) for some
\(a\in\mathfrak g\).
\item[(ii bis)]\SGARefTarget{sga3:xiii:proposition:4:4:item:ii-bis}
\(\mathfrak d\) is nilpotent, and there exists \(a\in\mathfrak d\)
such that \(\operatorname{ad}(a)_{\mathfrak g/\mathfrak d}\) is
injective.
\item[(iii)]\SGARefTarget{sga3:xiii:proposition:4:4:item:iii}
\(\mathfrak d\) is nilpotent and equal to its own normalizer.
\end{enumerate}
Then
\[
  \SGARefLink{sga3:xiii:proposition:4:4:item:i}{\text{\rm(i)}}
  \Longleftrightarrow
  \text{\rm\SGARefLink{sga3:xiii:proposition:4:4:item:i-bis}{(i bis)}}
  \Longrightarrow
  \SGARefLink{sga3:xiii:proposition:4:4:item:ii}{\text{\rm(ii)}}
  \Longleftrightarrow
  \text{\rm\SGARefLink{sga3:xiii:proposition:4:4:item:ii-bis}{(ii bis)}}
  \Longleftrightarrow
  \SGARefLink{sga3:xiii:proposition:4:4:item:iii}{\text{\rm(iii)}}.
\]
We shall see in \SGARefLink{sga3:xiii:corollary:5:7:part:a}{5.7 a)} that if \(\mathfrak g\) is the Lie algebra of
a smooth algebraic group, then all the preceding conditions are
equivalent.
\end{proposition}

The equivalence of \SGARefLink{sga3:xiii:proposition:4:4:item:i}{(i)} and \SGARefLink{sga3:xiii:proposition:4:4:item:i-bis}{(i bis)} follows immediately from \SGARefLink{sga3:xiii:corollary:4:3}{4.3}, and
these conditions plainly imply \SGARefLink{sga3:xiii:proposition:4:4:item:ii}{(ii)}. The equivalence of \SGARefLink{sga3:xiii:proposition:4:4:item:ii}{(ii)} and
\SGARefLink{sga3:xiii:proposition:4:4:item:ii-bis}{(ii bis)} is also immediate, as is \SGARefLink{sga3:xiii:proposition:4:4:item:ii-bis}{(ii bis)} \(\Rightarrow\) \SGARefLink{sga3:xiii:proposition:4:4:item:iii}{(iii)}
(cf. \SGARefLink{sga3:xiii:proposition:4:1}{4.1}). It remains to prove \SGARefLink{sga3:xiii:proposition:4:4:item:iii}{(iii)} \(\Rightarrow\) \SGARefLink{sga3:xiii:proposition:4:4:item:ii-bis}{(ii bis)}. This
is the only implication that uses the infinitude of \(k\), and it
follows at once from \SGARefLink{sga3:xiii:lemma:4:5}{the next lemma}.

\begin{lemma}
\SGARefTarget{sga3:xiii:lemma:4:5}
Let \(\mathfrak d\) be a nilpotent Lie algebra over an infinite field
\(k\), acting on a finite-dimensional vector space \(V\). Suppose
that the endomorphism \(u(x)\) is not injective for every
\(x\in\mathfrak d\). Then there is a nonzero element \(v\in V\)
annihilated by \(\mathfrak d\).
\end{lemma}

We may suppose that \(k\) is algebraically closed and that
\(\mathfrak d\) is finite-dimensional. Then \(V\) is known to be a
direct sum of finitely many nonzero stable vector subspaces \(V_i\),
\(1\leq i\leq n\), such that, for every \(i\) and every
\(x\in\mathfrak d\), the restriction \(u(x)|_{V_i}\) has a single
eigenvalue \(\lambda_i(x)\) (cf. Bourbaki, \emph{Groupes et Algèbres
de Lie}, Chap.~I, \(\S\)4, Exercise~22). Let \(c_i(x)\) be the
constant term of the characteristic polynomial of
\(u(x)|_{V_i}\), so that
\[
  \lambda_i(x)=0\quad\Longleftrightarrow\quad c_i(x)=0.
\]
Each \(c_i\) is a polynomial function on \(\mathfrak d\), and the
hypothesis says that \(\mathfrak d\) is the union of their zero sets.
Thus one of the \(c_i\) is identically zero. Replacing \(V\) by
\(V_i\), we are reduced to the case in which all \(u(x)\),
\(x\in\mathfrak d\), are nilpotent. Engel's theorem (Bourbaki,
loc.\ cit., Th.~1) then implies that some nonzero \(v\in V\) is
annihilated by \(\mathfrak d\).
\par\hfill\textsc{q.e.d.}

It is easy to see, still assuming that \(k\) is infinite, that
conditions \SGARefLink{sga3:xiii:proposition:4:4:item:i}{(i)} and \SGARefLink{sga3:xiii:proposition:4:4:item:i-bis}{(i bis)} of \SGARefLink{sga3:xiii:proposition:4:4}{4.4} are invariant under every extension
of the base field. When they hold, \(\mathfrak d\) is called a
\emph{Cartan subalgebra} of \(\mathfrak g\). In general, with \(k\)
not necessarily infinite, \(\mathfrak d\) is called a Cartan
subalgebra of \(\mathfrak g\) if it becomes one after one, and hence
every, extension \(k\to k'\) with \(k'\) infinite. It follows in
particular that \(\mathfrak d\) is nilpotent and equal to its own
normalizer.

\begin{proposition}
\SGARefTarget{sga3:xiii:proposition:4:6}
\SGARefTarget{sga3:xiii:proposition:4:6:item:a}
\textup{a)} Let \(a\in\mathfrak g\). If \(a\) is regular, then it is
contained in one and only one Cartan subalgebra of \(\mathfrak g\).
We shall see the converse in \SGARefLink{sga3:xiii:theorem:6:1:part:d}{6.1 d)}, when \(k\) is algebraically
closed and \(\mathfrak g\) is the Lie algebra of a smooth algebraic
group.

\SGARefTarget{sga3:xiii:proposition:4:6:item:b}
\textup{b)} Let \(\mathfrak d\) be a Cartan subalgebra of
\(\mathfrak g\), and let \(a\in\mathfrak d\). Then \(a\) is regular
in \(\mathfrak g\) if and only if
\(\operatorname{ad}(a)_{\mathfrak g/\mathfrak d}\) is injective.
\end{proposition}

For \SGARefLink{sga3:xiii:proposition:4:6:item:a}{a)}, if \(a\) is regular, then
\(\operatorname{Nil}(a,\mathfrak g)\) is a Cartan subalgebra of
\(\mathfrak g\), since this is true after an infinite extension
\(k'\) of \(k\). It follows at once from \SGARefLink{sga3:xiii:corollary:4:3}{4.3} that every Cartan
subalgebra of \(\mathfrak g\) containing \(a\) is equal to this one.
For \SGARefLink{sga3:xiii:proposition:4:6:item:b}{b)}, observe that the nullity of
\(\operatorname{ad}(a)_{\mathfrak g}\) is the sum of the nullities of
\(\operatorname{ad}(a)_{\mathfrak d}\) and
\(\operatorname{ad}(a)_{\mathfrak g/\mathfrak d}\). Since the first
is \(r\), their sum is \(r\) if and only if
\(\operatorname{ad}(a)_{\mathfrak g/\mathfrak d}\) is injective.
\par\hfill\textsc{q.e.d.}

\begin{corollary}
\SGARefTarget{sga3:xiii:corollary:4:7}
Let \(a\) be a regular element of \(\mathfrak g\), let
\(\mathfrak d\) be a Cartan subalgebra of \(\mathfrak g\) containing
\(a\), let \(A\) be a \(k\)-algebra, and let
\(\mathfrak g_A\) and \(\mathfrak d_A\subset\mathfrak g_A\) be the
\(A\)-Lie algebras obtained from \(\mathfrak g\) and \(\mathfrak d\)
by base change; write \(a_A\) for the image of \(a\) in
\(\mathfrak g_A\). Let \(u\) be an automorphism of
\(\mathfrak g_A\). Then
\[
  u(\mathfrak d_A)=\mathfrak d_A
\]
if and only if
\[
  u(a_A)\in\mathfrak d_A.
\]
\end{corollary}

The condition is plainly necessary; let us prove that it is
sufficient. Suppose it holds. Then
\[
  \mathfrak d'=u^{-1}(\mathfrak d_A)
\]
is a Lie subalgebra containing \(a_A\), and every \(b\in\mathfrak d'\)
is such that \(\operatorname{ad}(b)_{\mathfrak d'}\) is
nilpotent.\footnote{The source prints
\(\mathfrak d'=u(\mathfrak d_A)\) and
\(\operatorname{ad}(b)_{\mathfrak d}\). Under the hypothesis
\(u(a_A)\in\mathfrak d_A\), the inverse image
\(u^{-1}(\mathfrak d_A)\) is the subalgebra that contains \(a_A\),
and the nilpotence assertion concerns its restriction
\(\operatorname{ad}(b)_{\mathfrak d'}\).}
Indeed, \(\mathfrak d'\) is isomorphic to \(\mathfrak d_A\), which
has this property, as follows immediately from the definition of
``nilpotent'' in Bourbaki, \emph{Groupes et Algèbres de Lie},
Chap.~I, \(\S\)4, Def.~1. Taking \(b=a_A\), we see that
\(\operatorname{Nil}(b,\mathfrak g_A)\) contains \(\mathfrak d'\).
On the other hand,
\[
  \operatorname{Nil}(b,\mathfrak g_A)=\mathfrak d_A.
\]
Moreover, \(\mathfrak d'\) is locally a direct summand of
\(\mathfrak g_A\), since \(\mathfrak d\) is, and hence is locally a
direct summand of \(\mathfrak d_A\). It is a projective module of the
same rank \(r\) as \(\mathfrak d_A\), and therefore equals
\(\mathfrak d_A\).
\par\hfill\textsc{q.e.d.}

\begin{proposition}
\SGARefTarget{sga3:xiii:proposition:4:8}
Let \(\mathfrak h\) be a Lie subalgebra of \(\mathfrak g\).

\SGARefTarget{sga3:xiii:proposition:4:8:item:a}
\textup{a)} If \(k\) is infinite, the following conditions are
equivalent:
\begin{enumerate}
\item[(i)]\SGARefTarget{sga3:xiii:proposition:4:8:item:a:i}
\(\mathfrak h\) contains a Cartan subalgebra \(\mathfrak d\) of
\(\mathfrak g\).
\item[(ii)]\SGARefTarget{sga3:xiii:proposition:4:8:item:a:ii}
\(\mathfrak h\) contains a regular element \(a\) of \(\mathfrak g\),
and an element \(b\) such that
\(\operatorname{ad}(b)_{\mathfrak g/\mathfrak h}\) is injective.
\item[(iii)]\SGARefTarget{sga3:xiii:proposition:4:8:item:a:iii}
\(\mathfrak h\) has the same nilpotent rank as \(\mathfrak g\), and
contains a regular element of \(\mathfrak g\).
\end{enumerate}
These conditions are invariant under extension of the base field
\(k\).

\SGARefTarget{sga3:xiii:proposition:4:8:item:b}
\textup{b)} Suppose that these conditions hold after a suitable
infinite extension \(k'\) of \(k\). For \(a\in\mathfrak h\), the
element \(a\) is regular in \(\mathfrak g\) if and only if it is
regular in \(\mathfrak h\) and
\(\operatorname{ad}(a)_{\mathfrak g/\mathfrak h}\) is injective;
equivalently, if and only if
\[
  \operatorname{Nil}(a,\mathfrak h)
  =
  \operatorname{Nil}(a,\mathfrak g)
\]
and this common Lie algebra is a Cartan subalgebra of
\(\mathfrak h\).

\SGARefTarget{sga3:xiii:proposition:4:8:item:c}
\textup{c)} Under the hypotheses of \SGARefLink{sga3:xiii:proposition:4:8:item:b}{b)}, let \(\mathfrak d\) be a Lie
subalgebra of \(\mathfrak h\). In order that \(\mathfrak d\) be a
Cartan subalgebra of \(\mathfrak h\), it is sufficient that it be a
Cartan subalgebra of \(\mathfrak g\). We shall see in \SGARefLink{sga3:xiii:corollary:5:8}{5.8} that \SGARefLink{sga3:xiii:proposition:4:8:item:c}{this
condition} is also necessary when \(\mathfrak g\) is the Lie algebra
of a smooth algebraic group \(G\), and \(\mathfrak h\) is the Lie
algebra of a smooth algebraic subgroup \(H\) of \(G\).
\end{proposition}

Conditions \SGARefLink{sga3:xiii:proposition:4:8:item:a:ii}{(ii)} and \SGARefLink{sga3:xiii:proposition:4:8:item:a:iii}{(iii)} of \SGARefLink{sga3:xiii:proposition:4:8:item:a}{a)} are visibly invariant under extension
of the base field \(k\), assumed infinite. In \SGARefLink{sga3:xiii:proposition:4:8:item:b}{b)} and \SGARefLink{sga3:xiii:proposition:4:8:item:c}{c)}, we may
suppose that \(k\) is infinite, and shall do so. If \(\mathfrak h\)
contains the Cartan subalgebra
\[
  \mathfrak d=\operatorname{Nil}(a,\mathfrak g),
\]
then \(a\) is a regular element of \(\mathfrak g\) such that
\(\operatorname{ad}(a)_{\mathfrak g/\mathfrak h}\) is injective.
Thus \SGARefLink{sga3:xiii:proposition:4:8:item:a:i}{(i)} \(\Rightarrow\) \SGARefLink{sga3:xiii:proposition:4:8:item:a:ii}{(ii)}. Conversely, if \SGARefLink{sga3:xiii:proposition:4:8:item:a:ii}{(ii)} holds, then a
``sufficiently general'' element \(a\) of \(\mathfrak h\)
simultaneously satisfies both conditions in \SGARefLink{sga3:xiii:proposition:4:8:item:a:ii}{(ii)}. Hence
\(\operatorname{Nil}(a,\mathfrak g)\) is a Cartan subalgebra of
\(\mathfrak g\) contained in \(\mathfrak h\), proving \SGARefLink{sga3:xiii:proposition:4:8:item:a:i}{(i)}.
Conditions \SGARefLink{sga3:xiii:proposition:4:8:item:a:i}{(i)} and \SGARefLink{sga3:xiii:proposition:4:8:item:a:ii}{(ii)} are therefore equivalent.

Suppose they hold, and let \(a\) vary in \(\mathfrak h\). Then
\SGARefTarget{sga3:xiii:equation:4:8:x}
\[
\tag{x}
  \operatorname{rank}_k\operatorname{Nil}(a,\mathfrak g)
  =
  \operatorname{rank}_k\operatorname{Nil}(a,\mathfrak h)
  +
  \operatorname{rank}_k
  \operatorname{Nil}\bigl(
    \operatorname{ad}(a)_{\mathfrak g/\mathfrak h}
  \bigr).
\]
The two terms on the right are respectively at least \(r'\), the
nilpotent rank of \(\mathfrak h\), and at least \(0\); moreover, both
equalities are attained for a sufficiently general
element of \(\mathfrak h\). We also have
\[
  \operatorname{rank}_k\operatorname{Nil}(a,\mathfrak g)\geq r,
\]
where \(r\) is the nilpotent rank of \(\mathfrak g\). Equality is
attained for a sufficiently general element of \(\mathfrak h\), and
holds if and only if \(a\) is regular in \(\mathfrak g\).

It follows that \(r=r'\), and that \(a\) is regular if and only if
the two terms on the right-hand side of \(\SGARefLink{sga3:xiii:equation:4:8:x}{\mathrm{(x)}}\) are
respectively \(r'\) and \(0\); that is, if and only if \(a\) is
regular in \(\mathfrak h\) and
\(\operatorname{ad}(a)_{\mathfrak g/\mathfrak h}\) is injective.
This proves \SGARefLink{sga3:xiii:proposition:4:8:item:b}{b)}, and \SGARefLink{sga3:xiii:proposition:4:8:item:c}{c)} follows immediately by taking an element
\(a\in\mathfrak d\) regular in \(\mathfrak g\), so that
\(\operatorname{Nil}(a,\mathfrak g)=\mathfrak d\).

\SGARefLink{sga3:xiii:proposition:4:8:item:b}{The preceding result} also proves \SGARefLink{sga3:xiii:proposition:4:8:item:a:i}{(i)} \(\Rightarrow\) \SGARefLink{sga3:xiii:proposition:4:8:item:a:iii}{(iii)}. Finally,
\SGARefLink{sga3:xiii:proposition:4:8:item:a:iii}{(iii)} \(\Rightarrow\) \SGARefLink{sga3:xiii:proposition:4:8:item:a:i}{(i)}: under \SGARefLink{sga3:xiii:proposition:4:8:item:a:iii}{(iii)}, a sufficiently general
\(a\in\mathfrak h\) is regular both in \(\mathfrak h\) and in
\(\mathfrak g\). Thus
\[
  \operatorname{Nil}(a,\mathfrak h)
  \subset
  \operatorname{Nil}(a,\mathfrak g)
\]
are respectively Cartan subalgebras of \(\mathfrak h\) and
\(\mathfrak g\). Since they have the same rank over \(k\), they are
equal, proving \SGARefLink{sga3:xiii:proposition:4:8:item:a:i}{(i)}. This completes the proof of \SGARefLink{sga3:xiii:proposition:4:8}{4.8}.
